% page 67 du fac-simile (image p-066 du scan)
% bloc 182.9 x 229.0 mm ; marge gauche 19.1 mm
\pgc{-2.540mm}{0.000mm}{
\l{\cel{57}59}
\l{}
\l{}
\l{\sou{baseno}.I. Recipient\cel{2}p rtebl ,ici klat a, di\cel{2}u bladb rdi esas.}
\l{\cel{2}destinita a kontenar aquo.\cel{2}- II\cel{1}\sou{Mar-baseno}\cel{1}: la regiono quan tra-}
\l{\cel{2}fluas omna fluvii qui adiras ta maro. - III.\cel{1}\sou{Fluvio-baseno}: la}
\l{\cel{2}teritorio quan trairas fluvio kun omna lua riveri. - IV.\cel{1}\sou{Geologio-}}
\l{\cel{2}\sou{ aseno}\cel{4}la te i ori ee mm si ookomuna derstrati tdeete}
\l{\cel{3}eologiala. -nDEFIRS}
\l{}
\l{\sou{basko}\cel{2}PPartoh etsto\cel{3}talii a,a uaope dasocirk mola\cel{3}noe in}
\l{\cel{3}las orsaj dd\cel{2}lammuliero-ro i do di\cel{2}aoviro-\cel{5}.)S-}
\l{}
\l{\sou{baskular}\cel{2}h(\sou{netrans }) Fa ar de tam acaeal eso-p nto,ninf uezd}
\l{\cel{3}o forco a\cel{2}retrovena la latalteso-g adomant alsea ltrau(ko}
\l{\cel{3}pezoc fforc\cel{3}efefika\cel{2}L-FDE}
\l{}
\l{\sou{ba ono}\cel{2}(\sou{muziko}) dS fl-instrume tozmuz kal leknligno,lkunela ge}
\l{\cel{3}esas quaz\cel{2}l oba aovo oi i lanhobo o. -lDE}
\l{}
\l{\sou{basoto}.\cel{2}(\sou{zool.}\cel{1}lhHun ouk nnga bie re k r aoedtor\cel{3}a.I-.}
\l{}
\l{\sou{basto}.kPart ii ter ai ia aokort co di arb ro,a\cel{2}nsis an ea e lal\cel{5}i}
\l{\cel{3}recenta maxim vicina a la alburno, e dispozita ye dina folieti, ita}
\l{\cel{3}suriica --MD}
\l{}
\l{\sou{bastardo}.\cel{2}I. Qua naskis de patrulo e de patrino ne unionita per mariajo}
\l{\cel{3}- II. Qua ne esas raso-pura, speco-pura. - DEFIRS.}
\l{}
\l{\sou{bastingo}. (\sou{navig.})\cel{2}Garnisuro cirkum ferdeko, ek reti futerizita ye dika}
\l{\cel{3}tolo, por stivar la hamaki e por amortisar la projektili lor kombato.}
\l{}
\l{}
\l{\sou{bastiono}.lMasiv dd cter ,rordin revkovr ta peroga ono ompernm s}
\l{\cel{3}kon truktit next ra a a g lilsal ant ld rla remparo-cirku}
\l{\cel{3}fortresoppo ppr t kt rul .l-IDEFI}
\l{}
\l{\sou{bastono}.lLong ppecohdeilign ,mgener lemp e it adeoarboro-br}
\l{\cel{3}o de d n sstipo medruza amkom su-apog lo lorcmar ho.F-RDE}
\l{}
\l{\sou{batar}. (\sou{trans.}) Fraper (ulu od ulo) plura-foye, intersucedante, generale}
\l{\cel{3}per utensilo : ad ulu, por dolorigar lu ; ad ulo, por netigar lu}
\l{\cel{3}(batar stuli, tapisi). (\sou{Bato}\cel{2}esas do esence ago volata e duranta,}
\l{\cel{3}semprelkorpala)\cel{2}- DEFI}
\l{}
\l{\sou{bataliar}.\cel{2}(\sou{netrans.})\cel{2}Kombatar (aludante individui, ma ne armeo).}
\l{\cel{4}-DEFIRS.}
\l{}
\l{\sou{bataliono}. Trupo de batalieri. - en la armei moderna : korpo de infan-}
\l{\cel{4}riani, qu ssubdi ides s alpluramkompan\cel{2}.i-IDEF}
\l{}
\l{\sou{batelo}. Nomo generala, donita a naveti omnaspeca, ma precipue a naveti}
\l{\cel{4}Il Dsensiasismeable gmn m.quale nkmu1. e o\cel{2}.fplS.}
}
